% !TeX program = xelatex
% !TeX encoding = UTF-8
% Sinh tu du lieu that: nghien cuu MS-GARCH-X + EGARCH-X qua
% scripts/bao_cao/sinh_pr_a_pr_b_tex.py -- khong nhap tay. Phu luc S1-S8
% (ban minh hoa cu, B thap) da BO -- xem git history neu can, thay bang
% tong_hop_ket_qua.tex (gop S1-S8 vao mot bo bang duy nhat, khong trung
% lap voi nghien cuu MS-GARCH-X chinh thuc).
\documentclass[11pt,a4paper]{article}
\usepackage{fontspec}
\usepackage{unicode-math}
\IfFontExistsTF{Times New Roman}{\setmainfont{Times New Roman}}{\setmainfont{Liberation Serif}}
\IfFontExistsTF{Cambria Math}{\setmathfont{Cambria Math}}{\setmathfont{LibertinusMath-Regular.otf}}
\usepackage[a4paper,top=2.2cm,bottom=2.2cm,left=1.8cm,right=1.8cm]{geometry}
\usepackage{booktabs}
\usepackage{tabularx}
\usepackage{array}
\usepackage{threeparttable}
\usepackage{longtable}
\usepackage{float}
\usepackage{caption}
\usepackage{xcolor}
\usepackage[hidelinks]{hyperref}
\usepackage{microtype}
\emergencystretch=3em
\sloppy
\definecolor{Accent}{HTML}{2E5A88}
\definecolor{Warn}{HTML}{9B4B4B}
\captionsetup{font=small,labelfont=bf,skip=5pt}
\renewcommand{\tablename}{Bảng}
\setlength{\parindent}{0pt}
\setlength{\parskip}{0.4em}
\newcolumntype{C}{>{\centering\arraybackslash}X}
\newcolumntype{L}{>{\raggedright\arraybackslash}X}

\title{\vspace{-1.2cm}\textbf{Nghiên cứu mô phỏng khôi phục tham số --- Bước 2}\\[0.25em]
\large MS-GARCH-X và EGARCH-X Student-$t$}
\author{}
\date{}

\begin{document}
\maketitle
\vspace{-1.4cm}

\section*{Tổng quan}
Tài liệu này trình bày \textbf{toàn bộ phần nghiên cứu mô phỏng khôi phục tham số
chính thức} (Bước 2 trong thiết kế ba bước hiện hành của đề tài --- xem
\texttt{CLAUDE.md}, mục "Quyết định thiết kế đã chốt"), cộng thêm phần xác
nhận chéo với Bước 3:

\begin{itemize}
\item \textbf{MS-GARCH-X Student-$t$}, hai trạng thái (15 tham số),
  3 kịch bản $S_1,S_2,S_3$ $\times$ 3 phương pháp (MLE, VB mean-field, VB
  full-rank) $\times$ $B=100$ lần lặp, $T=2000$. \textbf{Đã hoàn tất 900/900
  lần lặp.}
\item \textbf{EGARCH-X Student-$t$}, một trạng thái (8 tham số),
  3 kịch bản $E_1,E_2,E_3$ $\times$ 3 phương pháp $\times$ $B=100$, $T=2000$.
  \textbf{Đã hoàn tất 900/900 lần lặp.}
\item \textbf{Bước 3} --- MCMC một lần duy nhất trên EGARCH-X Student-$t$ áp
  lên dữ liệu VN30F1M thật, \textbf{đã hội tụ} ($\hat R=1{,}0125$). Bảng
  $R_{sd}$ xác nhận trên dữ liệu thật đúng phát hiện "VB đánh giá thấp bất
  định ở $\omega,\beta$" mà hai nghiên cứu mô phỏng trên đã chỉ ra.
\end{itemize}

\textbf{Về bộ tám kịch bản minh họa cũ ($S_1$--$S_8$, $B=15$--$30$, trước khi
sửa lỗi sắp nhãn trạng thái):} tài liệu này không còn trình bày riêng bộ đó
nữa, để tránh hai bộ số khác nhau cho cùng tên kịch bản $S_1,S_2,S_3$ trong
một tài liệu. Bộ tám kịch bản gộp đúng cách (dùng $B=100$ cho $S_1$--$S_3$ từ
nghiên cứu MS-GARCH-X chính thức, $B=29$/$15$ cho $S_4$--$S_8$) nằm trong
\texttt{Bài báo/tong\_hop\_ket\_qua.tex}.

\textbf{Sửa lỗi quan trọng đã áp dụng cho cả hai nghiên cứu mô phỏng:}
\begin{itemize}
\item \textbf{Đối xứng hoán vị nhãn trạng thái} (chỉ ở MS-GARCH-X, vì EGARCH-X
  là mô hình một trạng thái không có đối xứng này): bộ ước lượng có thể trả về "trạng thái 1"/
  "trạng thái 2" theo thứ tự tùy ý. Đã sắp nhãn theo phương sai vô điều kiện
  $\omega_s/(1-\alpha_s-\beta_s)$ (trạng thái 1 = biến động thấp) trước khi tính
  Bias/RMSE/CP. \textbf{VB đảo nhãn gần 50\% lần lặp} --- gần đúng bằng tung đồng
  xu, một phát hiện độc lập đáng chú ý.
\item \textbf{Placeholder $p_{11}=p_{22}=1{,}0$}: giá trị thật để tính Bias/RMSE/CP
  phải lấy từ $P_0$ của kịch bản, không phải placeholder $\mathrm{eye}(2)$
  trong \texttt{scenario.params.P}. Lỗi này từng lọt qua \texttt{run\_scenario()}
  (đường tuần tự) lẫn \texttt{combine()} (đường song song, đã sinh ra toàn bộ
  900 lần lặp của nghiên cứu MS-GARCH-X) --- phát hiện và vá độc lập hai lần, lần sau (21/09/2026)
  đổi $p_{11}/p_{22}$ đáng kể (ví dụ $S_1$: MLE $p_{11}=0{,}983$ thay vì
  $0{,}963$ cũ; VB-FR coverage của $p_{11},p_{22}$ tăng từ gần $0$ lên
  $0{,}23$--$0{,}69$). Bảng dưới đây đã dùng số liệu sau lần vá thứ hai.
\item \textbf{Brier tính sai vì tính trước khi sắp nhãn}: đã chạy lại bộ lọc
  tiến với tham số đã lưu (không cần ước lượng lại) để tính Brier đúng --- giá
  trị trong caption mỗi bảng dưới đây đã là bản đúng.
\item \textbf{Coverage của MLE chưa từng được tính}: nay tính bằng khoảng Wald
  $\hat\theta\pm1{,}96\,\mathrm{SE}$ trên thang tự nhiên, SE từ Hessian quan
  sát. \textbf{Hạn chế cần biết}: Hessian của MS-GARCH-X chỉ xác định dương ở
  khoảng một nửa số lần lặp (ví dụ $S_1$: 50/100) --- cột CP của MLE dưới đây
  là \emph{coverage có điều kiện trên tập con Hessian hợp lệ}, không phải trên
  toàn bộ $B=100$ lần lặp.
\end{itemize}

\clearpage
\section{MS-GARCH-X Student-$t$, hai trạng thái}

Ba kịch bản: $S_1$ (mốc nền, hai trạng thái phân biệt rõ), $S_2$ (nhận dạng yếu,
hai trạng thái gần nhau --- cũng là nơi VB mean-field và VB full-rank khác nhau rõ
nhất vì mean-field bỏ tương quan hậu nghiệm), $S_3$ (đuôi rất dày, $\nu_0=4$ sát
giá trị ước lượng thật trên VN30F1M $\hat\nu\approx4{,}08$).

Cách trình bày bảng: giá trị thật $\theta_0$, trung bình ước lượng qua các lần
lặp (Mean $=\theta_0+$ Độ chệch), độ chệch, RMSE, độ lệch chuẩn Monte Carlo (SE,
suy ra từ $\mathrm{SE}=\sqrt{\max(\mathrm{RMSE}^2-\mathrm{Bias}^2,0)}$), và độ
phủ khoảng tin cậy 95\% (CP). Với VB, CP dùng phân vị hậu nghiệm; với MLE, CP
dùng khoảng Wald từ Hessian quan sát, chỉ tính trên các lần lặp có Hessian xác
định dương (khoảng một nửa số lần lặp đối với MS-GARCH-X --- xem phần Tổng
quan ở đầu tài liệu).

\subsection*{Kịch bản S1 --- Mốc nền --- hai trạng thái phân biệt rõ, $\nu_0=8$}
\begin{table}[H]\centering\normalsize
\caption{S1: toàn bộ tham số thật $\theta_0$ dùng để sinh dữ liệu}
\begin{tabularx}{\textwidth}{l C  l C}
\toprule
\textbf{Tham số} & \textbf{Giá trị} & \textbf{Tham số} & \textbf{Giá trị} \\
\midrule
$\mu_1$ & 0,000 & $\gamma_{1,1}$ & 0,050 \\
$\mu_2$ & 0,000 & $\gamma_{1,2}$ & 0,050 \\
$\omega_1$ & 0,050 & $\gamma_{2,1}$ & 0,050 \\
$\omega_2$ & 0,400 & $\gamma_{2,2}$ & 0,050 \\
$\alpha_1$ & 0,050 & $p_{11}$ & 0,980 \\
$\alpha_2$ & 0,100 & $p_{22}$ & 0,950 \\
$\beta_1$ & 0,850 & $\nu$ & 8,000 \\
$\beta_2$ & 0,800 &  &  \\
\bottomrule
\end{tabularx}
\vspace{0.2em}\raggedright\scriptsize \emph{Ghi chú:} $T=2000$, burn-in $=200$, $B=100$ lần lặp (mục tiêu và thực đạt đều bằng nhau, trừ vài lần lặp phân kỳ NaN đã dùng quy tắc thử lại điểm khởi tạo).
\end{table}

\begin{table}[H]\centering\footnotesize
\caption{Hợp lý cực đại (MLE) ($n=100$ lần lặp, Brier trung bình $=0,3969$)}
\begin{tabularx}{\textwidth}{l C C C C C C}
\toprule
\textbf{Tham số} & \textbf{Thật ($\theta_0$)} & \textbf{Mean} & \textbf{Độ chệch} & \textbf{RMSE} & \textbf{SE} & \textbf{CP 95\%} \\
\midrule
$\mu_1$ & 0,000 & $-$0,002 & $-$0,002 & 0,023 & 0,023 & 84\% \\
$\mu_2$ & 0,000 & 0,004 & 0,004 & 0,135 & 0,135 & 82\% \\
$\omega_1$ & 0,050 & 0,085 & 0,035 & 0,054 & 0,041 & 84\% \\
$\omega_2$ & 0,400 & 1,076 & 0,676 & 1,152 & 0,933 & 82\% \\
$\alpha_1$ & 0,050 & 0,084 & 0,034 & 0,058 & 0,047 & 88\% \\
$\alpha_2$ & 0,100 & 0,142 & 0,042 & 0,102 & 0,093 & 88\% \\
$\beta_1$ & 0,850 & 0,720 & $-$0,130 & 0,160 & 0,094 & 60\% \\
$\beta_2$ & 0,800 & 0,733 & $-$0,067 & 0,256 & 0,247 & 90\% \\
$\gamma_{1,1}$ & 0,050 & 0,060 & 0,010 & 0,027 & 0,026 & 86\% \\
$\gamma_{1,2}$ & 0,050 & 0,061 & 0,011 & 0,027 & 0,025 & 82\% \\
$\gamma_{2,1}$ & 0,050 & 0,066 & 0,016 & 0,266 & 0,265 & 88\% \\
$\gamma_{2,2}$ & 0,050 & 0,122 & 0,072 & 0,324 & 0,316 & 82\% \\
$p_{11}$ & 0,980 & 0,983 & 0,003 & 0,018 & 0,018 & 52\% \\
$p_{22}$ & 0,950 & 0,945 & $-$0,005 & 0,025 & 0,024 & 92\% \\
$\nu$ & 8,000 & 6,397 & $-$1,603 & 1,953 & 1,117 & 62\% \\
\bottomrule
\end{tabularx}
\end{table}

\begin{table}[H]\centering\footnotesize
\caption{Bayes biến phân mean-field ($n=100$ lần lặp, Brier trung bình $=0,4947$)}
\begin{tabularx}{\textwidth}{l C C C C C C}
\toprule
\textbf{Tham số} & \textbf{Thật ($\theta_0$)} & \textbf{Mean} & \textbf{Độ chệch} & \textbf{RMSE} & \textbf{SE} & \textbf{CP 95\%} \\
\midrule
$\mu_1$ & 0,000 & $-$0,008 & $-$0,008 & 0,099 & 0,099 & 64\% \\
$\mu_2$ & 0,000 & 0,000 & 0,000 & 0,112 & 0,112 & 59\% \\
$\omega_1$ & 0,050 & 0,186 & 0,136 & 0,154 & 0,072 & 0\% \\
$\omega_2$ & 0,400 & 0,232 & $-$0,168 & 0,190 & 0,089 & 16\% \\
$\alpha_1$ & 0,050 & 0,344 & 0,294 & 0,308 & 0,091 & 0\% \\
$\alpha_2$ & 0,100 & 0,365 & 0,265 & 0,278 & 0,086 & 0\% \\
$\beta_1$ & 0,850 & 0,673 & $-$0,177 & 0,186 & 0,057 & 0\% \\
$\beta_2$ & 0,800 & 0,655 & $-$0,145 & 0,154 & 0,051 & 0\% \\
$\gamma_{1,1}$ & 0,050 & 0,062 & 0,012 & 0,071 & 0,070 & 90\% \\
$\gamma_{1,2}$ & 0,050 & 0,078 & 0,028 & 0,093 & 0,088 & 79\% \\
$\gamma_{2,1}$ & 0,050 & 0,043 & $-$0,007 & 0,090 & 0,089 & 87\% \\
$\gamma_{2,2}$ & 0,050 & 0,047 & $-$0,003 & 0,088 & 0,087 & 86\% \\
$p_{11}$ & 0,980 & 0,810 & $-$0,170 & 0,184 & 0,071 & 0\% \\
$p_{22}$ & 0,950 & 0,724 & $-$0,226 & 0,236 & 0,069 & 3\% \\
$\nu$ & 8,000 & 3,739 & $-$4,261 & 4,340 & 0,821 & 0\% \\
\bottomrule
\end{tabularx}
\end{table}

\begin{table}[H]\centering\footnotesize
\caption{Bayes biến phân full-rank ($n=100$ lần lặp, Brier trung bình $=0,4689$)}
\begin{tabularx}{\textwidth}{l C C C C C C}
\toprule
\textbf{Tham số} & \textbf{Thật ($\theta_0$)} & \textbf{Mean} & \textbf{Độ chệch} & \textbf{RMSE} & \textbf{SE} & \textbf{CP 95\%} \\
\midrule
$\mu_1$ & 0,000 & $-$0,007 & $-$0,007 & 0,067 & 0,067 & 86\% \\
$\mu_2$ & 0,000 & 0,036 & 0,036 & 0,160 & 0,156 & 75\% \\
$\omega_1$ & 0,050 & 0,224 & 0,174 & 0,209 & 0,117 & 1\% \\
$\omega_2$ & 0,400 & 0,542 & 0,142 & 0,338 & 0,307 & 50\% \\
$\alpha_1$ & 0,050 & 0,290 & 0,240 & 0,256 & 0,091 & 1\% \\
$\alpha_2$ & 0,100 & 0,374 & 0,274 & 0,291 & 0,099 & 18\% \\
$\beta_1$ & 0,850 & 0,575 & $-$0,275 & 0,299 & 0,117 & 0\% \\
$\beta_2$ & 0,800 & 0,605 & $-$0,195 & 0,213 & 0,085 & 33\% \\
$\gamma_{1,1}$ & 0,050 & 0,047 & $-$0,003 & 0,065 & 0,065 & 91\% \\
$\gamma_{1,2}$ & 0,050 & 0,051 & 0,001 & 0,061 & 0,061 & 92\% \\
$\gamma_{2,1}$ & 0,050 & 0,058 & 0,008 & 0,121 & 0,121 & 95\% \\
$\gamma_{2,2}$ & 0,050 & 0,069 & 0,019 & 0,156 & 0,155 & 91\% \\
$p_{11}$ & 0,980 & 0,906 & $-$0,074 & 0,109 & 0,080 & 40\% \\
$p_{22}$ & 0,950 & 0,763 & $-$0,187 & 0,219 & 0,113 & 23\% \\
$\nu$ & 8,000 & 4,483 & $-$3,517 & 3,677 & 1,071 & 8\% \\
\bottomrule
\end{tabularx}
\end{table}

\emph{Đảo nhãn trạng thái trước khi gộp (đã sửa):} MLE 23/100, VB-MF 48/100, VB-FR 49/100 lần lặp trả về nhãn trạng thái ngược quy ước (đã hoán vị trước khi tính Bias/RMSE/CP ở trên).
\bigskip

\subsection*{Kịch bản S2 --- Nhận dạng yếu --- hai trạng thái gần nhau, $\nu_0=8$}
\begin{table}[H]\centering\normalsize
\caption{S2: toàn bộ tham số thật $\theta_0$ dùng để sinh dữ liệu}
\begin{tabularx}{\textwidth}{l C  l C}
\toprule
\textbf{Tham số} & \textbf{Giá trị} & \textbf{Tham số} & \textbf{Giá trị} \\
\midrule
$\mu_1$ & 0,000 & $\gamma_{1,1}$ & 0,050 \\
$\mu_2$ & 0,000 & $\gamma_{1,2}$ & 0,050 \\
$\omega_1$ & 0,080 & $\gamma_{2,1}$ & 0,050 \\
$\omega_2$ & 0,140 & $\gamma_{2,2}$ & 0,050 \\
$\alpha_1$ & 0,060 & $p_{11}$ & 0,950 \\
$\alpha_2$ & 0,080 & $p_{22}$ & 0,900 \\
$\beta_1$ & 0,850 & $\nu$ & 8,000 \\
$\beta_2$ & 0,830 &  &  \\
\bottomrule
\end{tabularx}
\vspace{0.2em}\raggedright\scriptsize \emph{Ghi chú:} $T=2000$, burn-in $=200$, $B=100$ lần lặp (mục tiêu và thực đạt đều bằng nhau, trừ vài lần lặp phân kỳ NaN đã dùng quy tắc thử lại điểm khởi tạo).
\end{table}

\begin{table}[H]\centering\footnotesize
\caption{Hợp lý cực đại (MLE) ($n=100$ lần lặp, Brier trung bình $=0,6611$)}
\begin{tabularx}{\textwidth}{l C C C C C C}
\toprule
\textbf{Tham số} & \textbf{Thật ($\theta_0$)} & \textbf{Mean} & \textbf{Độ chệch} & \textbf{RMSE} & \textbf{SE} & \textbf{CP 95\%} \\
\midrule
$\mu_1$ & 0,000 & $-$0,011 & $-$0,011 & 0,178 & 0,178 & 80\% \\
$\mu_2$ & 0,000 & 0,062 & 0,062 & 0,343 & 0,337 & 63\% \\
$\omega_1$ & 0,080 & 0,108 & 0,028 & 0,086 & 0,081 & 82\% \\
$\omega_2$ & 0,140 & 0,646 & 0,506 & 0,994 & 0,855 & 69\% \\
$\alpha_1$ & 0,060 & 0,051 & $-$0,009 & 0,045 & 0,044 & 76\% \\
$\alpha_2$ & 0,080 & 0,087 & 0,007 & 0,104 & 0,103 & 78\% \\
$\beta_1$ & 0,850 & 0,804 & $-$0,046 & 0,147 & 0,139 & 84\% \\
$\beta_2$ & 0,830 & 0,752 & $-$0,078 & 0,312 & 0,302 & 76\% \\
$\gamma_{1,1}$ & 0,050 & 0,061 & 0,011 & 0,068 & 0,068 & 78\% \\
$\gamma_{1,2}$ & 0,050 & 0,055 & 0,005 & 0,114 & 0,114 & 73\% \\
$\gamma_{2,1}$ & 0,050 & 0,061 & 0,011 & 0,364 & 0,364 & 76\% \\
$\gamma_{2,2}$ & 0,050 & 0,049 & $-$0,001 & 0,378 & 0,378 & 73\% \\
$p_{11}$ & 0,950 & 0,960 & 0,010 & 0,076 & 0,076 & 45\% \\
$p_{22}$ & 0,900 & 0,878 & $-$0,022 & 0,123 & 0,121 & 59\% \\
$\nu$ & 8,000 & 8,667 & 0,667 & 1,718 & 1,583 & 100\% \\
\bottomrule
\end{tabularx}
\end{table}

\begin{table}[H]\centering\footnotesize
\caption{Bayes biến phân mean-field ($n=100$ lần lặp, Brier trung bình $=0,4976$)}
\begin{tabularx}{\textwidth}{l C C C C C C}
\toprule
\textbf{Tham số} & \textbf{Thật ($\theta_0$)} & \textbf{Mean} & \textbf{Độ chệch} & \textbf{RMSE} & \textbf{SE} & \textbf{CP 95\%} \\
\midrule
$\mu_1$ & 0,000 & 0,000 & 0,000 & 0,080 & 0,080 & 81\% \\
$\mu_2$ & 0,000 & $-$0,002 & $-$0,002 & 0,101 & 0,101 & 81\% \\
$\omega_1$ & 0,080 & 0,388 & 0,308 & 0,317 & 0,072 & 0\% \\
$\omega_2$ & 0,140 & 0,415 & 0,275 & 0,287 & 0,080 & 0\% \\
$\alpha_1$ & 0,060 & 0,144 & 0,084 & 0,088 & 0,028 & 17\% \\
$\alpha_2$ & 0,080 & 0,154 & 0,074 & 0,082 & 0,036 & 43\% \\
$\beta_1$ & 0,850 & 0,499 & $-$0,351 & 0,358 & 0,069 & 0\% \\
$\beta_2$ & 0,830 & 0,510 & $-$0,320 & 0,326 & 0,062 & 0\% \\
$\gamma_{1,1}$ & 0,050 & 0,051 & 0,001 & 0,055 & 0,055 & 98\% \\
$\gamma_{1,2}$ & 0,050 & 0,063 & 0,013 & 0,052 & 0,050 & 95\% \\
$\gamma_{2,1}$ & 0,050 & 0,047 & $-$0,003 & 0,049 & 0,049 & 96\% \\
$\gamma_{2,2}$ & 0,050 & 0,041 & $-$0,009 & 0,058 & 0,057 & 93\% \\
$p_{11}$ & 0,950 & 0,791 & $-$0,159 & 0,163 & 0,037 & 5\% \\
$p_{22}$ & 0,900 & 0,767 & $-$0,133 & 0,140 & 0,044 & 58\% \\
$\nu$ & 8,000 & 7,218 & $-$0,782 & 1,396 & 1,157 & 75\% \\
\bottomrule
\end{tabularx}
\end{table}

\begin{table}[H]\centering\footnotesize
\caption{Bayes biến phân full-rank ($n=100$ lần lặp, Brier trung bình $=0,4894$)}
\begin{tabularx}{\textwidth}{l C C C C C C}
\toprule
\textbf{Tham số} & \textbf{Thật ($\theta_0$)} & \textbf{Mean} & \textbf{Độ chệch} & \textbf{RMSE} & \textbf{SE} & \textbf{CP 95\%} \\
\midrule
$\mu_1$ & 0,000 & $-$0,012 & $-$0,012 & 0,065 & 0,064 & 89\% \\
$\mu_2$ & 0,000 & 0,011 & 0,011 & 0,160 & 0,159 & 75\% \\
$\omega_1$ & 0,080 & 0,397 & 0,317 & 0,329 & 0,089 & 0\% \\
$\omega_2$ & 0,140 & 0,618 & 0,478 & 0,520 & 0,205 & 0\% \\
$\alpha_1$ & 0,060 & 0,131 & 0,071 & 0,077 & 0,030 & 45\% \\
$\alpha_2$ & 0,080 & 0,210 & 0,130 & 0,143 & 0,060 & 69\% \\
$\beta_1$ & 0,850 & 0,462 & $-$0,388 & 0,396 & 0,079 & 0\% \\
$\beta_2$ & 0,830 & 0,464 & $-$0,366 & 0,374 & 0,074 & 1\% \\
$\gamma_{1,1}$ & 0,050 & 0,051 & 0,001 & 0,047 & 0,047 & 100\% \\
$\gamma_{1,2}$ & 0,050 & 0,066 & 0,016 & 0,051 & 0,049 & 94\% \\
$\gamma_{2,1}$ & 0,050 & 0,067 & 0,017 & 0,107 & 0,106 & 96\% \\
$\gamma_{2,2}$ & 0,050 & 0,048 & $-$0,002 & 0,131 & 0,131 & 94\% \\
$p_{11}$ & 0,950 & 0,885 & $-$0,065 & 0,086 & 0,056 & 69\% \\
$p_{22}$ & 0,900 & 0,714 & $-$0,186 & 0,195 & 0,061 & 25\% \\
$\nu$ & 8,000 & 7,704 & $-$0,296 & 1,402 & 1,371 & 94\% \\
\bottomrule
\end{tabularx}
\end{table}

\emph{Đảo nhãn trạng thái trước khi gộp (đã sửa):} MLE 30/100, VB-MF 52/100, VB-FR 47/100 lần lặp trả về nhãn trạng thái ngược quy ước (đã hoán vị trước khi tính Bias/RMSE/CP ở trên).
\bigskip

\subsection*{Kịch bản S3 --- Đuôi rất dày --- $\nu_0=4$, sát biên tồn tại phương sai}
\begin{table}[H]\centering\normalsize
\caption{S3: toàn bộ tham số thật $\theta_0$ dùng để sinh dữ liệu}
\begin{tabularx}{\textwidth}{l C  l C}
\toprule
\textbf{Tham số} & \textbf{Giá trị} & \textbf{Tham số} & \textbf{Giá trị} \\
\midrule
$\mu_1$ & 0,000 & $\gamma_{1,1}$ & 0,050 \\
$\mu_2$ & 0,000 & $\gamma_{1,2}$ & 0,050 \\
$\omega_1$ & 0,050 & $\gamma_{2,1}$ & 0,050 \\
$\omega_2$ & 0,400 & $\gamma_{2,2}$ & 0,050 \\
$\alpha_1$ & 0,050 & $p_{11}$ & 0,980 \\
$\alpha_2$ & 0,100 & $p_{22}$ & 0,950 \\
$\beta_1$ & 0,850 & $\nu$ & 4,000 \\
$\beta_2$ & 0,800 &  &  \\
\bottomrule
\end{tabularx}
\vspace{0.2em}\raggedright\scriptsize \emph{Ghi chú:} $T=2000$, burn-in $=200$, $B=100$ lần lặp (mục tiêu và thực đạt đều bằng nhau, trừ vài lần lặp phân kỳ NaN đã dùng quy tắc thử lại điểm khởi tạo).
\end{table}

\begin{table}[H]\centering\footnotesize
\caption{Hợp lý cực đại (MLE) ($n=100$ lần lặp, Brier trung bình $=0,3992$)}
\begin{tabularx}{\textwidth}{l C C C C C C}
\toprule
\textbf{Tham số} & \textbf{Thật ($\theta_0$)} & \textbf{Mean} & \textbf{Độ chệch} & \textbf{RMSE} & \textbf{SE} & \textbf{CP 95\%} \\
\midrule
$\mu_1$ & 0,000 & $-$0,001 & $-$0,001 & 0,020 & 0,020 & 83\% \\
$\mu_2$ & 0,000 & $-$0,008 & $-$0,008 & 0,111 & 0,111 & 81\% \\
$\omega_1$ & 0,050 & 0,090 & 0,040 & 0,057 & 0,040 & 83\% \\
$\omega_2$ & 0,400 & 1,096 & 0,696 & 1,434 & 1,253 & 74\% \\
$\alpha_1$ & 0,050 & 0,071 & 0,021 & 0,048 & 0,043 & 83\% \\
$\alpha_2$ & 0,100 & 0,134 & 0,034 & 0,113 & 0,108 & 90\% \\
$\beta_1$ & 0,850 & 0,682 & $-$0,168 & 0,201 & 0,110 & 52\% \\
$\beta_2$ & 0,800 & 0,776 & $-$0,024 & 0,260 & 0,259 & 74\% \\
$\gamma_{1,1}$ & 0,050 & 0,059 & 0,009 & 0,029 & 0,027 & 86\% \\
$\gamma_{1,2}$ & 0,050 & 0,063 & 0,013 & 0,029 & 0,026 & 88\% \\
$\gamma_{2,1}$ & 0,050 & 0,142 & 0,092 & 0,387 & 0,376 & 83\% \\
$\gamma_{2,2}$ & 0,050 & 0,147 & 0,097 & 0,321 & 0,306 & 90\% \\
$p_{11}$ & 0,980 & 0,979 & $-$0,001 & 0,018 & 0,018 & 62\% \\
$p_{22}$ & 0,950 & 0,945 & $-$0,005 & 0,038 & 0,038 & 81\% \\
$\nu$ & 4,000 & 3,936 & $-$0,064 & 0,537 & 0,534 & 76\% \\
\bottomrule
\end{tabularx}
\end{table}

\begin{table}[H]\centering\footnotesize
\caption{Bayes biến phân mean-field ($n=100$ lần lặp, Brier trung bình $=0,4763$)}
\begin{tabularx}{\textwidth}{l C C C C C C}
\toprule
\textbf{Tham số} & \textbf{Thật ($\theta_0$)} & \textbf{Mean} & \textbf{Độ chệch} & \textbf{RMSE} & \textbf{SE} & \textbf{CP 95\%} \\
\midrule
$\mu_1$ & 0,000 & 0,008 & 0,008 & 0,079 & 0,079 & 64\% \\
$\mu_2$ & 0,000 & $-$0,006 & $-$0,006 & 0,103 & 0,103 & 60\% \\
$\omega_1$ & 0,050 & 0,273 & 0,223 & 0,234 & 0,071 & 0\% \\
$\omega_2$ & 0,400 & 0,335 & $-$0,065 & 0,119 & 0,100 & 45\% \\
$\alpha_1$ & 0,050 & 0,541 & 0,491 & 0,501 & 0,099 & 0\% \\
$\alpha_2$ & 0,100 & 0,541 & 0,441 & 0,450 & 0,089 & 0\% \\
$\beta_1$ & 0,850 & 0,716 & $-$0,134 & 0,143 & 0,050 & 0\% \\
$\beta_2$ & 0,800 & 0,691 & $-$0,109 & 0,123 & 0,057 & 11\% \\
$\gamma_{1,1}$ & 0,050 & 0,120 & 0,070 & 0,113 & 0,089 & 81\% \\
$\gamma_{1,2}$ & 0,050 & 0,138 & 0,088 & 0,136 & 0,104 & 72\% \\
$\gamma_{2,1}$ & 0,050 & 0,099 & 0,049 & 0,110 & 0,098 & 87\% \\
$\gamma_{2,2}$ & 0,050 & 0,110 & 0,060 & 0,115 & 0,098 & 89\% \\
$p_{11}$ & 0,980 & 0,799 & $-$0,181 & 0,189 & 0,053 & 0\% \\
$p_{22}$ & 0,950 & 0,731 & $-$0,219 & 0,225 & 0,053 & 0\% \\
$\nu$ & 4,000 & 2,404 & $-$1,596 & 1,602 & 0,137 & 0\% \\
\bottomrule
\end{tabularx}
\end{table}

\begin{table}[H]\centering\footnotesize
\caption{Bayes biến phân full-rank ($n=100$ lần lặp, Brier trung bình $=0,4431$)}
\begin{tabularx}{\textwidth}{l C C C C C C}
\toprule
\textbf{Tham số} & \textbf{Thật ($\theta_0$)} & \textbf{Mean} & \textbf{Độ chệch} & \textbf{RMSE} & \textbf{SE} & \textbf{CP 95\%} \\
\midrule
$\mu_1$ & 0,000 & $-$0,006 & $-$0,006 & 0,049 & 0,049 & 93\% \\
$\mu_2$ & 0,000 & $-$0,016 & $-$0,016 & 0,125 & 0,124 & 82\% \\
$\omega_1$ & 0,050 & 0,281 & 0,231 & 0,264 & 0,129 & 0\% \\
$\omega_2$ & 0,400 & 0,673 & 0,273 & 0,418 & 0,317 & 51\% \\
$\alpha_1$ & 0,050 & 0,361 & 0,311 & 0,332 & 0,116 & 3\% \\
$\alpha_2$ & 0,100 & 0,431 & 0,331 & 0,345 & 0,096 & 19\% \\
$\beta_1$ & 0,850 & 0,565 & $-$0,285 & 0,313 & 0,129 & 1\% \\
$\beta_2$ & 0,800 & 0,613 & $-$0,187 & 0,212 & 0,099 & 30\% \\
$\gamma_{1,1}$ & 0,050 & 0,077 & 0,027 & 0,080 & 0,075 & 86\% \\
$\gamma_{1,2}$ & 0,050 & 0,088 & 0,038 & 0,080 & 0,070 & 84\% \\
$\gamma_{2,1}$ & 0,050 & 0,080 & 0,030 & 0,145 & 0,142 & 92\% \\
$\gamma_{2,2}$ & 0,050 & 0,095 & 0,045 & 0,187 & 0,181 & 89\% \\
$p_{11}$ & 0,980 & 0,903 & $-$0,077 & 0,101 & 0,066 & 24\% \\
$p_{22}$ & 0,950 & 0,779 & $-$0,171 & 0,199 & 0,101 & 25\% \\
$\nu$ & 4,000 & 2,795 & $-$1,205 & 1,240 & 0,290 & 4\% \\
\bottomrule
\end{tabularx}
\end{table}

\emph{Đảo nhãn trạng thái trước khi gộp (đã sửa):} MLE 28/100, VB-MF 48/100, VB-FR 53/100 lần lặp trả về nhãn trạng thái ngược quy ước (đã hoán vị trước khi tính Bias/RMSE/CP ở trên).
\bigskip

\textbf{Đọc kết quả MS-GARCH-X (tóm tắt):} sau khi vá lỗi $p_{11}/p_{22}$ (21/09/2026),
bức tranh coverage của hai tham số này khác hẳn báo cáo trước --- \textbf{không
còn đúng} để nói "coverage của $p_{11},p_{22}$ bằng $0{,}00$ ở mọi kịch bản,
mọi phương pháp Bayes" như từng ghi nhận nhầm. Số liệu ở trên: MLE (khoảng
Wald, có điều kiện trên Hessian hợp lệ) đạt $45$--$92\%$ tùy tham số/kịch bản;
VB-FR đạt $23$--$69\%$, rõ ràng khá hơn VB-MF ($0$--$58\%$, phần lớn dưới
$10\%$ trừ một trường hợp ở $S_2$). Tiên nghiệm Beta$(8;2)$ (trung bình
$0{,}80$) kéo VB-MF về gần giá trị đó vẫn là cơ chế chính lý giải vì sao
VB-MF còn yếu ở $p_{ss}$, nhưng không còn đúng để nói VB \emph{luôn luôn}
thất bại ở tham số này --- VB-FR (nắm được tương quan hậu nghiệm) cải thiện
rõ rệt so với VB-MF. Ở $\omega_2$ (tham số nhận dạng yếu khác), MLE vẫn có
đuôi chạy trốn (cực đại vượt $5{,}0$ ở $S_1$) trong khi VB tránh được nhờ
tiên nghiệm LogNormal --- đây là nơi VB thắng RMSE rõ nhất.

\clearpage
\section{EGARCH-X Student-$t$, một trạng thái}

Ba kịch bản: $E_1$ (đòn bẩy rõ $\theta_0=-0{,}10$), $E_2$ (không đòn bẩy
$\theta_0=0$ --- \textbf{kịch bản quan trọng nhất}, kiểm tra mô hình có bịa ra
hiệu ứng đòn bẩy khi dữ liệu vốn không có không), $E_3$ (đuôi dày $\nu_0=4$, có
đòn bẩy). Mô hình một trạng thái nên \textbf{không có đối xứng hoán vị nhãn} --
đây là lý do EGARCH-X không gặp vấn đề đảo nhãn như MS-GARCH-X.

\subsection*{Kịch bản E1 --- Đòn bẩy rõ, $\theta_0=-0{,}10$}
\begin{table}[H]\centering\normalsize
\caption{E1: toàn bộ tham số thật $\theta_0$ dùng để sinh dữ liệu}
\begin{tabularx}{\textwidth}{l C  l C}
\toprule
\textbf{Tham số} & \textbf{Giá trị} & \textbf{Tham số} & \textbf{Giá trị} \\
\midrule
$\mu$ & 0,050 & $\beta$ & 0,950 \\
$\omega$ & 0,040 & $\gamma_1$ & 0,050 \\
$\alpha$ & 0,150 & $\gamma_2$ & 0,050 \\
$\theta$ & $-$0,100 & $\nu$ & 8,000 \\
\bottomrule
\end{tabularx}
\vspace{0.2em}\raggedright\scriptsize \emph{Ghi chú:} $T=2000$, $B=100$ lần lặp, mô hình một trạng thái (không có đối xứng hoán vị nhãn).
\end{table}

\begin{table}[H]\centering\footnotesize
\caption{Hợp lý cực đại (MLE) ($n=100$ lần lặp)}
\begin{tabularx}{\textwidth}{l C C C C C C}
\toprule
\textbf{Tham số} & \textbf{Thật ($\theta_0$)} & \textbf{Mean} & \textbf{Độ chệch} & \textbf{RMSE} & \textbf{SE} & \textbf{CP 95\%} \\
\midrule
$\mu$ & 0,050 & 0,047 & $-$0,003 & 0,031 & 0,031 & 97\% \\
$\omega$ & 0,040 & 0,043 & 0,003 & 0,013 & 0,012 & 93\% \\
$\alpha$ & 0,150 & 0,148 & $-$0,002 & 0,028 & 0,028 & 94\% \\
$\theta$ & $-$0,100 & $-$0,103 & $-$0,003 & 0,018 & 0,018 & 95\% \\
$\beta$ & 0,950 & 0,946 & $-$0,004 & 0,014 & 0,014 & 95\% \\
$\gamma_1$ & 0,050 & 0,051 & 0,001 & 0,018 & 0,018 & 92\% \\
$\gamma_2$ & 0,050 & 0,054 & 0,004 & 0,018 & 0,017 & 94\% \\
$\nu$ & 8,000 & 8,570 & 0,570 & 1,648 & 1,546 & 97\% \\
\bottomrule
\end{tabularx}
\end{table}

\begin{table}[H]\centering\footnotesize
\caption{Bayes biến phân mean-field ($n=100$ lần lặp)}
\begin{tabularx}{\textwidth}{l C C C C C C}
\toprule
\textbf{Tham số} & \textbf{Thật ($\theta_0$)} & \textbf{Mean} & \textbf{Độ chệch} & \textbf{RMSE} & \textbf{SE} & \textbf{CP 95\%} \\
\midrule
$\mu$ & 0,050 & 0,047 & $-$0,003 & 0,031 & 0,031 & 94\% \\
$\omega$ & 0,040 & 0,063 & 0,023 & 0,031 & 0,021 & 32\% \\
$\alpha$ & 0,150 & 0,191 & 0,041 & 0,048 & 0,025 & 68\% \\
$\theta$ & $-$0,100 & $-$0,111 & $-$0,011 & 0,023 & 0,020 & 89\% \\
$\beta$ & 0,950 & 0,922 & $-$0,028 & 0,035 & 0,020 & 12\% \\
$\gamma_1$ & 0,050 & 0,055 & 0,005 & 0,022 & 0,022 & 91\% \\
$\gamma_2$ & 0,050 & 0,059 & 0,009 & 0,022 & 0,020 & 91\% \\
$\nu$ & 8,000 & 8,632 & 0,632 & 1,570 & 1,437 & 95\% \\
\bottomrule
\end{tabularx}
\end{table}

\begin{table}[H]\centering\footnotesize
\caption{Bayes biến phân full-rank ($n=100$ lần lặp)}
\begin{tabularx}{\textwidth}{l C C C C C C}
\toprule
\textbf{Tham số} & \textbf{Thật ($\theta_0$)} & \textbf{Mean} & \textbf{Độ chệch} & \textbf{RMSE} & \textbf{SE} & \textbf{CP 95\%} \\
\midrule
$\mu$ & 0,050 & 0,047 & $-$0,003 & 0,032 & 0,032 & 96\% \\
$\omega$ & 0,040 & 0,076 & 0,036 & 0,060 & 0,049 & 66\% \\
$\alpha$ & 0,150 & 0,203 & 0,053 & 0,075 & 0,052 & 67\% \\
$\theta$ & $-$0,100 & $-$0,114 & $-$0,014 & 0,026 & 0,022 & 95\% \\
$\beta$ & 0,950 & 0,908 & $-$0,042 & 0,068 & 0,053 & 54\% \\
$\gamma_1$ & 0,050 & 0,055 & 0,005 & 0,023 & 0,023 & 87\% \\
$\gamma_2$ & 0,050 & 0,059 & 0,009 & 0,023 & 0,022 & 91\% \\
$\nu$ & 8,000 & 8,548 & 0,548 & 1,601 & 1,504 & 96\% \\
\bottomrule
\end{tabularx}
\end{table}

\bigskip

\subsection*{Kịch bản E2 --- Không đòn bẩy, $\theta_0=0$ --- kịch bản quan trọng nhất}
\begin{table}[H]\centering\normalsize
\caption{E2: toàn bộ tham số thật $\theta_0$ dùng để sinh dữ liệu}
\begin{tabularx}{\textwidth}{l C  l C}
\toprule
\textbf{Tham số} & \textbf{Giá trị} & \textbf{Tham số} & \textbf{Giá trị} \\
\midrule
$\mu$ & 0,050 & $\beta$ & 0,950 \\
$\omega$ & 0,040 & $\gamma_1$ & 0,050 \\
$\alpha$ & 0,150 & $\gamma_2$ & 0,050 \\
$\theta$ & 0,000 & $\nu$ & 8,000 \\
\bottomrule
\end{tabularx}
\vspace{0.2em}\raggedright\scriptsize \emph{Ghi chú:} $T=2000$, $B=100$ lần lặp, mô hình một trạng thái (không có đối xứng hoán vị nhãn).
\end{table}

\begin{table}[H]\centering\footnotesize
\caption{Hợp lý cực đại (MLE) ($n=100$ lần lặp)}
\begin{tabularx}{\textwidth}{l C C C C C C}
\toprule
\textbf{Tham số} & \textbf{Thật ($\theta_0$)} & \textbf{Mean} & \textbf{Độ chệch} & \textbf{RMSE} & \textbf{SE} & \textbf{CP 95\%} \\
\midrule
$\mu$ & 0,050 & 0,049 & $-$0,001 & 0,031 & 0,031 & 97\% \\
$\omega$ & 0,040 & 0,044 & 0,004 & 0,016 & 0,016 & 95\% \\
$\alpha$ & 0,150 & 0,149 & $-$0,001 & 0,029 & 0,029 & 95\% \\
$\theta$ & 0,000 & $-$0,001 & $-$0,001 & 0,017 & 0,017 & 97\% \\
$\beta$ & 0,950 & 0,944 & $-$0,006 & 0,019 & 0,018 & 95\% \\
$\gamma_1$ & 0,050 & 0,051 & 0,001 & 0,018 & 0,018 & 93\% \\
$\gamma_2$ & 0,050 & 0,056 & 0,006 & 0,019 & 0,018 & 96\% \\
$\nu$ & 8,000 & 8,572 & 0,572 & 1,643 & 1,540 & 97\% \\
\bottomrule
\end{tabularx}
\end{table}

\begin{table}[H]\centering\footnotesize
\caption{Bayes biến phân mean-field ($n=100$ lần lặp)}
\begin{tabularx}{\textwidth}{l C C C C C C}
\toprule
\textbf{Tham số} & \textbf{Thật ($\theta_0$)} & \textbf{Mean} & \textbf{Độ chệch} & \textbf{RMSE} & \textbf{SE} & \textbf{CP 95\%} \\
\midrule
$\mu$ & 0,050 & 0,048 & $-$0,002 & 0,033 & 0,033 & 96\% \\
$\omega$ & 0,040 & 0,101 & 0,061 & 0,080 & 0,051 & 8\% \\
$\alpha$ & 0,150 & 0,213 & 0,063 & 0,069 & 0,029 & 45\% \\
$\theta$ & 0,000 & 0,000 & 0,000 & 0,021 & 0,021 & 96\% \\
$\beta$ & 0,950 & 0,875 & $-$0,075 & 0,095 & 0,058 & 2\% \\
$\gamma_1$ & 0,050 & 0,059 & 0,009 & 0,025 & 0,024 & 92\% \\
$\gamma_2$ & 0,050 & 0,065 & 0,015 & 0,028 & 0,024 & 89\% \\
$\nu$ & 8,000 & 8,497 & 0,497 & 1,502 & 1,417 & 94\% \\
\bottomrule
\end{tabularx}
\end{table}

\begin{table}[H]\centering\footnotesize
\caption{Bayes biến phân full-rank ($n=100$ lần lặp)}
\begin{tabularx}{\textwidth}{l C C C C C C}
\toprule
\textbf{Tham số} & \textbf{Thật ($\theta_0$)} & \textbf{Mean} & \textbf{Độ chệch} & \textbf{RMSE} & \textbf{SE} & \textbf{CP 95\%} \\
\midrule
$\mu$ & 0,050 & 0,049 & $-$0,001 & 0,033 & 0,033 & 97\% \\
$\omega$ & 0,040 & 0,109 & 0,069 & 0,092 & 0,061 & 35\% \\
$\alpha$ & 0,150 & 0,219 & 0,069 & 0,083 & 0,046 & 50\% \\
$\theta$ & 0,000 & $-$0,003 & $-$0,003 & 0,023 & 0,023 & 98\% \\
$\beta$ & 0,950 & 0,866 & $-$0,084 & 0,112 & 0,073 & 23\% \\
$\gamma_1$ & 0,050 & 0,059 & 0,009 & 0,026 & 0,024 & 90\% \\
$\gamma_2$ & 0,050 & 0,064 & 0,014 & 0,028 & 0,025 & 90\% \\
$\nu$ & 8,000 & 8,478 & 0,478 & 1,546 & 1,471 & 97\% \\
\bottomrule
\end{tabularx}
\end{table}

\bigskip

\subsection*{Kịch bản E3 --- Đuôi dày, $\nu_0=4$, có đòn bẩy}
\begin{table}[H]\centering\normalsize
\caption{E3: toàn bộ tham số thật $\theta_0$ dùng để sinh dữ liệu}
\begin{tabularx}{\textwidth}{l C  l C}
\toprule
\textbf{Tham số} & \textbf{Giá trị} & \textbf{Tham số} & \textbf{Giá trị} \\
\midrule
$\mu$ & 0,050 & $\beta$ & 0,950 \\
$\omega$ & 0,040 & $\gamma_1$ & 0,050 \\
$\alpha$ & 0,150 & $\gamma_2$ & 0,050 \\
$\theta$ & $-$0,100 & $\nu$ & 4,000 \\
\bottomrule
\end{tabularx}
\vspace{0.2em}\raggedright\scriptsize \emph{Ghi chú:} $T=2000$, $B=100$ lần lặp, mô hình một trạng thái (không có đối xứng hoán vị nhãn).
\end{table}

\begin{table}[H]\centering\footnotesize
\caption{Hợp lý cực đại (MLE) ($n=100$ lần lặp)}
\begin{tabularx}{\textwidth}{l C C C C C C}
\toprule
\textbf{Tham số} & \textbf{Thật ($\theta_0$)} & \textbf{Mean} & \textbf{Độ chệch} & \textbf{RMSE} & \textbf{SE} & \textbf{CP 95\%} \\
\midrule
$\mu$ & 0,050 & 0,048 & $-$0,002 & 0,026 & 0,026 & 93\% \\
$\omega$ & 0,040 & 0,043 & 0,003 & 0,013 & 0,012 & 95\% \\
$\alpha$ & 0,150 & 0,146 & $-$0,004 & 0,032 & 0,032 & 89\% \\
$\theta$ & $-$0,100 & $-$0,103 & $-$0,003 & 0,022 & 0,021 & 95\% \\
$\beta$ & 0,950 & 0,946 & $-$0,004 & 0,015 & 0,014 & 92\% \\
$\gamma_1$ & 0,050 & 0,050 & 0,000 & 0,021 & 0,021 & 92\% \\
$\gamma_2$ & 0,050 & 0,054 & 0,004 & 0,022 & 0,021 & 89\% \\
$\nu$ & 4,000 & 4,121 & 0,121 & 0,435 & 0,418 & 98\% \\
\bottomrule
\end{tabularx}
\end{table}

\begin{table}[H]\centering\footnotesize
\caption{Bayes biến phân mean-field ($n=100$ lần lặp)}
\begin{tabularx}{\textwidth}{l C C C C C C}
\toprule
\textbf{Tham số} & \textbf{Thật ($\theta_0$)} & \textbf{Mean} & \textbf{Độ chệch} & \textbf{RMSE} & \textbf{SE} & \textbf{CP 95\%} \\
\midrule
$\mu$ & 0,050 & 0,047 & $-$0,003 & 0,026 & 0,026 & 94\% \\
$\omega$ & 0,040 & 0,064 & 0,024 & 0,033 & 0,022 & 29\% \\
$\alpha$ & 0,150 & 0,189 & 0,039 & 0,048 & 0,028 & 75\% \\
$\theta$ & $-$0,100 & $-$0,113 & $-$0,013 & 0,025 & 0,022 & 91\% \\
$\beta$ & 0,950 & 0,919 & $-$0,031 & 0,038 & 0,021 & 12\% \\
$\gamma_1$ & 0,050 & 0,055 & 0,005 & 0,027 & 0,026 & 89\% \\
$\gamma_2$ & 0,050 & 0,059 & 0,009 & 0,026 & 0,025 & 89\% \\
$\nu$ & 4,000 & 4,188 & 0,188 & 0,464 & 0,424 & 83\% \\
\bottomrule
\end{tabularx}
\end{table}

\begin{table}[H]\centering\footnotesize
\caption{Bayes biến phân full-rank ($n=100$ lần lặp)}
\begin{tabularx}{\textwidth}{l C C C C C C}
\toprule
\textbf{Tham số} & \textbf{Thật ($\theta_0$)} & \textbf{Mean} & \textbf{Độ chệch} & \textbf{RMSE} & \textbf{SE} & \textbf{CP 95\%} \\
\midrule
$\mu$ & 0,050 & 0,047 & $-$0,003 & 0,027 & 0,027 & 94\% \\
$\omega$ & 0,040 & 0,076 & 0,036 & 0,054 & 0,040 & 69\% \\
$\alpha$ & 0,150 & 0,201 & 0,051 & 0,069 & 0,046 & 71\% \\
$\theta$ & $-$0,100 & $-$0,117 & $-$0,017 & 0,029 & 0,023 & 96\% \\
$\beta$ & 0,950 & 0,905 & $-$0,045 & 0,071 & 0,055 & 49\% \\
$\gamma_1$ & 0,050 & 0,055 & 0,005 & 0,026 & 0,026 & 90\% \\
$\gamma_2$ & 0,050 & 0,060 & 0,010 & 0,029 & 0,027 & 90\% \\
$\nu$ & 4,000 & 4,131 & 0,131 & 0,466 & 0,447 & 92\% \\
\bottomrule
\end{tabularx}
\end{table}

\bigskip

\textbf{Đọc kết quả EGARCH-X (tóm tắt):} bộ ước lượng khôi phục đòn bẩy đúng theo CẢ
HAI chiều --- $E_1$: 100/100 lần lặp cho $\hat\theta$ âm; $E_2$: $\hat\theta$
trung bình gần 0 ($-0{,}001$ đến $-0{,}003$, không bịa ra đòn bẩy giả); $E_3$:
100/100 âm, đúng dấu ngay cả khi đuôi rất dày. $\nu$ ở $E_3$ khôi phục gần như
hoàn hảo (MLE $4{,}131$ so với thật $4{,}000$), hóa giải lo ngại mô hình đánh
giá thấp độ dày đuôi đúng ở vùng quan trọng nhất của VN30F1M. \textbf{Coverage
của $\theta$ nay có cả cột MLE} (bổ sung 21/09/2026, Hessian một trạng thái
luôn xác định dương nên không cần điều kiện hóa như MS-GARCH-X): $95$--$98\%$ ở cả
ba phương pháp, cả ba kịch bản --- đạt chuẩn $0{,}95$ đồng đều, khác hẳn MS-GARCH-X.

\clearpage
\section{Bước 3 --- xác nhận chéo giữa mô phỏng và dữ liệu thật}

Hai nghiên cứu mô phỏng ở trên (MS-GARCH-X và EGARCH-X) kiểm định bộ ước
lượng trên dữ liệu \emph{mô phỏng}, nơi biết trước tham số thật. Bước 3 chạy
MCMC (NUTS) một lần duy nhất trên \textbf{EGARCH-X Student-$t$} --- mô hình
thắng cuộc của Bước 1 --- áp lên \textbf{dữ liệu VN30F1M thật}, cho một phép
kiểm định độc lập thứ hai của đúng những phát hiện đã nêu ở trên.

\textbf{MCMC hội tụ, khác hẳn lần thất bại trên MS-GARCH-X:} 4 chuỗi $\times$
(1.000 warmup + 400 mẫu), khởi tạo phân tán quanh nghiệm MLE, 8,1 giờ.
$\hat R$ lớn nhất $1{,}0125$, ESS nhỏ nhất $1.340/1.600$ mẫu (84\%), \textbf{0
phân kỳ} --- so với $\hat R=2{,}84$, ESS$\approx 5$, 64 giờ của lần MCMC thất
bại trước đây trên MS-GARCH-X. Đây là bằng chứng dứt khoát rằng thất bại cũ
đến từ \textbf{đối xứng hoán vị nhãn của cấu trúc chuyển trạng thái}, không
phải từ MCMC: cùng thư viện, cùng thuật toán, cùng dữ liệu --- bỏ cấu trúc
chuyển trạng thái thì hội tụ bình thường.

\begin{table}[H]\centering\small
\caption{Kết quả trên VN30F1M, bốn phương pháp cùng một hàm hợp lý}
\begin{tabularx}{\textwidth}{l C C C C L}
\toprule
\textbf{Tham số} & \textbf{MLE} & \textbf{VB-MF} & \textbf{VB-FR} & \textbf{MCMC} & \textbf{KTC 95\% (MCMC)} \\
\midrule
$\theta$ & $-$0{,}1452 & $-$0{,}1407 & $-$0{,}1241 & $-$0{,}1399 & $[-0{,}2125;\,-0{,}0710]$ \\
$\beta$ & 0{,}6832 & 0{,}6264 & 0{,}6390 & 0{,}6393 & $[0{,}4912;\,0{,}7562]$ \\
$\nu$ & 4{,}2910 & 4{,}3897 & 4{,}3647 & 4{,}3644 & $[3{,}5544;\,5{,}4315]$ \\
$\omega$ & 0{,}1599 & 0{,}1941 & 0{,}1897 & 0{,}1870 & $[0{,}1093;\,0{,}2948]$ \\
$\alpha$ & 0{,}0721 & 0{,}1081 & 0{,}1066 & 0{,}1121 & $[0{,}0268;\,0{,}2110]$ \\
\bottomrule
\end{tabularx}
\end{table}

\textbf{Hiệu ứng đòn bẩy của VN30F1M được xác lập thống kê}: khoảng tin cậy
95\% của $\theta$ loại trừ 0. Kết hợp với kịch bản $E_2$ của nghiên cứu EGARCH-X ($\theta_0=0$: mô hình
không bịa ra đòn bẩy khi dữ liệu vốn không có), kết luận này được kiểm chứng
\textbf{cả hai chiều} --- một lần trên dữ liệu mô phỏng có $\theta_0$ biết
trước, một lần trên dữ liệu thật không biết $\theta$ thật.

\subsection*{$R_{sd}$ --- tỷ số độ lệch chuẩn VB/MCMC, xác nhận trực tiếp phát hiện của hai nghiên cứu mô phỏng}

Coi MCMC là chuẩn tham chiếu (đã hội tụ), tỷ số $R_{sd}=\mathrm{sd}_{VB}/\mathrm{sd}_{MCMC}$
đo mức VB đánh giá thấp bất định hậu nghiệm --- $R_{sd}<1$ nghĩa là khoảng tin
cậy của VB hẹp hơn thực tế.

\begin{table}[H]\centering\small
\caption{$R_{sd}$ trên EGARCH-X Student-\text{$t$}, dữ liệu VN30F1M thật}
\begin{tabularx}{\textwidth}{l C C C C C C}
\toprule
 & $\mu$ & $\omega$ & $\alpha$ & $\theta$ & $\beta$ & $\nu$ \\
\midrule
VB-MF & 0{,}95 & \textbf{0{,}37} & 0{,}88 & 0{,}92 & \textbf{0{,}29} & 0{,}63 \\
VB-FR & 0{,}94 & \textbf{0{,}75} & 1{,}01 & 1{,}00 & \textbf{0{,}75} & 0{,}97 \\
\bottomrule
\end{tabularx}
\end{table}

Mean-field đánh giá thấp bất định của $\omega$ và $\beta$ khoảng \textbf{ba
lần} --- đúng hai tham số nằm trên sống núi cong $\omega/(1-\beta)$ mà cả
nghiên cứu MS-GARCH-X lẫn EGARCH-X đã chỉ ra là nơi coverage của VB tệ nhất.
Full-rank thu hẹp khoảng cách nhưng vẫn hẹp hơn MCMC khoảng 25\% ở đúng hai
tham số đó; ở $\mu,\theta,\alpha$ thì cả hai họ VB gần khớp MCMC. \textbf{Kết
quả mô phỏng (hai nghiên cứu ở trên) và kết quả dữ liệu thật (Bước 3) nói
cùng một điều, bằng hai con đường hoàn toàn độc lập} --- đây là dạng bằng
chứng mạnh nhất có thể đưa ra cho luận điểm về sống núi cong, vì mô phỏng
biết trước tham số thật còn dữ liệu thật thì không.

\end{document}
